% Part: first-order-logic
% Chapter: syntax-and-semantics
% Section: Structures

\documentclass[../../../include/open-logic-section]{subfiles}

\begin{document}

\olfileid{fol}{syn}{str}

\olsection{प्रथम-क्रम भाषांसाठीच्या \printtoken{P}{structure}}

\begin{explain}
प्रथम-क्रम भाषांना स्वतःचा कोणताही \emph{अर्थ लावलेला नसतो:}
!!{constant}s, !!{function}s आणि !!{predicate}s यांना कोणताही विशिष्ट
अर्थ जोडलेला नसतो. एखादी \article{structure} \emph{!!{structure}}
विनिर्दिष्ट करून अर्थ दिले जातात. ती \emph{प्रांत}, म्हणजे
!!{constant}s ज्या वस्तू निर्देशित करतात, !!{function}s ज्यांवर क्रिया
करतात आणि संख्यापक ज्या वस्तूंवर फिरतात त्या वस्तू विनिर्दिष्ट करते.
याशिवाय, कोणते !!{constant}s कोणत्या वस्तू निर्देशित करतात,
!!a{function} वस्तूंचे वस्तूंमध्ये प्रतिचित्रण कसे करते आणि !!{predicate}s
कोणत्या वस्तूंना लागू होतात हे ती विनिर्दिष्ट करते. !!^{structure}s या
तर्कशास्त्रातील \emph{चिन्हार्थविषयक} संकल्पनांचा---उदाहरणार्थ, तार्किक
निष्पन्नता, वैधता आणि पूर्ततायोग्यता या संकल्पनांचा---आधार आहेत. साहित्यात त्यांना
निरनिराळ्या ठिकाणी ``रचना,'' ``अर्थनिर्धारण'' किंवा ``प्रतिमान'' असे
म्हटले जाते.
\end{explain}

\begin{defn}[!!^{structure}s]
प्रथम-क्रम तर्कशास्त्राची भाषा~$\Lang{L}$ हिच्यासाठीची
\Article{structure} \emph{!!{structure}}~$\Struct M$ पुढील घटकांनी बनते:
\begin{enumerate}
\item \emph{प्रांत:} अरिक्त संच~$\Domain M$.
\item \emph{!!{constant}s चे अर्थनिर्धारण:} $\Lang{L}$ मधील प्रत्येक
  !!{constant}~$c$ साठी !!a{element}~$\Assign{c}{M} \in \Domain M$.
\item \emph{!!{predicate}s चे अर्थनिर्धारण:} $\Lang{L}$ मधील प्रत्येक
  $n$-स्थानी !!{predicate}~$R$ साठी ($\eq$ सोडून) $n$-स्थानी संबंध
  $\Assign{R}{M} \subseteq \Domain{M}^n$.
\item \emph{!!{function}s चे अर्थनिर्धारण:} $\Lang{L}$ मधील प्रत्येक
  $n$-स्थानी !!{function}~$f$ साठी $n$-स्थानी फलन
  $\Assign{f}{M} \colon \Domain{M}^n \to \Domain{M}$.
\end{enumerate}
\end{defn}

\begin{ex}
अंकगणिताच्या भाषेसाठीच्या !!^a{structure}~$\Struct M$ मध्ये एक संच,
!!{constant}~$\Obj 0$ चे अर्थनिर्धारण म्हणून $\Domain M$ मधील
घटक~$\Assign{\Obj 0}{M}$, एकस्थानी फलन
$\Assign{\Obj \prime}{M} \colon \Domain{M} \to \Domain M$,
$\Domain M^2 \to \Domain M$ अशी दोन द्विस्थानी फलने
$\Assign{\Obj +}{M}$ आणि $\Assign{\Obj \times}{M}$, तसेच द्विस्थानी
संबंध~$\Assign{\Obj <}{M} \subseteq \Domain{M}^2$ असतात.

अशा रचनेचे एक स्वाभाविक उदाहरण पुढीलप्रमाणे आहे:
\begin{enumerate}
\item $\Domain N = \Nat$
\item $\Assign{\Obj 0}{N} = 0$
\item प्रत्येक $n \in \Nat$ साठी $\Assign{\Obj \prime}{N}(n) = n + 1$
\item प्रत्येक $n, m \in \Nat$ साठी $\Assign{\Obj +}{N}(n, m) = n + m$
\item प्रत्येक $n, m \in \Nat$ साठी $\Assign{\Obj \times}{N}(n, m) = n\cdot m$
\item $\Assign{\Obj <}{N} = \Setabs{\tuple{n, m}}{n \in \Nat, m \in
  \Nat, n < m}$
\end{enumerate}
अशा रीतीने परिभाषित केलेल्या $\Lang L_A$ साठीच्या रचनेला~$\Struct N$
\emph{अंकगणिताचे मानक प्रतिमान} म्हणतात, कारण ती~$\Lang L_A$ मधील
अतार्किक चिन्हांचा अर्थ आपण अपेक्षा करू त्याच रीतीने लावते.

तरीही, $\Lang L_A$ साठी इतर अनेक !!{structure}s शक्य आहेत. उदाहरणार्थ,
प्रांत म्हणून~$\Nat$ ऐवजी पूर्णांकांचा संच~$\Int$ घेऊन $\Obj 0$,
$\Obj \prime$, $\Obj +$, $\Obj \times$, $\Obj <$ यांची अर्थनिर्धारणे
त्यानुसार परिभाषित करता येतील. पण संख्यांशी दुरान्वयानेही संबंध नसलेल्या
$\Lang L_A$ साठीच्या रचनाही आपण परिभाषित करू शकतो.
\end{ex}

\begin{ex}
संचसिद्धान्ताची भाषा~$\Lang L_Z$ हिच्यासाठीची रचना~$\Struct M$ हिला
फक्त एक संच आणि एकच द्विस्थानी संबंध आवश्यक असतात. म्हणून तांत्रिकदृष्ट्या,
उदाहरणार्थ, माणसांचा संच आणि ``$x$ हा~$y$ पेक्षा वयाने मोठा आहे'' हा संबंध
$\Lang L_Z$ साठी !!a{structure} म्हणून वापरता येईल; त्याचप्रमाणे~$\Nat$
आणि $n, m \in \Nat$ साठीचा $n \ge m$ हाही वापरता येईल.
%READERNOTE{OLFOL-010: गोठवलेल्या इंग्रजी मजकुरातील ``single-two place relation'' हे जोडचिन्ह चुकीचे आहे; व्याख्येला आवश्यक असलेला एकच द्विस्थानी संबंध मराठीत स्पष्ट केला आहे.}

$\Lang L_Z$ साठीची एक विशेष रोचक !!{structure}, जिच्या प्रांतातील
!!{element}s प्रत्यक्षात संच आहेत आणि जिच्यात $\Obj \in$ चे अर्थनिर्धारण
प्रत्यक्षात ``$x$ हा~$y$ चा !!a{element} आहे'' हा संबंध आहे, ती म्हणजे
\emph{आनुवंशिकदृष्ट्या सांत संचांची} !!{structure}~$\Struct{{HF}}$:
\begin{enumerate}
\item $\Domain{{{HF}}} = \emptyset \cup \Pow{\emptyset} \cup
  \Pow{\Pow{\emptyset}} \cup \Pow{\Pow{\Pow{\emptyset}}} \cup \dots$;
\item $\Assign{\Obj \in}{{{HF}}} = \Setabs{\tuple{x, y}}{x, y \in
  \Domain{{{HF}}}, x \in y}$.
\end{enumerate}
\end{ex}

\begin{digress}
!!a{structure} म्हणून काय गणले जाईल याविषयी आपण केलेल्या अटींचा आपल्या
तर्कशास्त्रावर परिणाम होतो. उदाहरणार्थ, प्रांत अरिक्त ठेवण्याची निवड
संख्यकीकृत वाक्यांच्या पूर्तिची (किंवा सत्यतेची) नेहमीची मांडणी गृहीत
धरल्यास $\lexists[x][(!A(x) \lor \lnot !A(x))]$ वैध---म्हणजे तार्किक
सत्य---असल्याची खात्री देते. तसेच, सर्व !!{constant}s नी प्रांतातील
वस्तूच निर्देशित केली पाहिजे ही अट अस्तित्ववाची सामान्यीकरण हा
निर्दोष अनुमान-रूपबंध असल्याची खात्री देते: $!A(a)$, म्हणून
$\lexists[x][!A(x)]$. नावांना प्रांताबाहेरील वस्तू निर्देशित करण्याची
किंवा काहीच निर्देशित न करण्याची मुभा दिली, तर आपण \emph{मुक्त
तर्कशास्त्राच्या} दिशेने जाऊ; त्यात अस्तित्ववाची सामान्यीकरणासाठी एक
अतिरिक्त आधारवाक्य लागते: $!A(a)$ आणि $\lexists[x][\eq[x][a]]$, म्हणून
$\lexists[x][!A(x)]$.
\end{digress}

\end{document}
